% ══════════════════════════════════════════════════════════════════
%  Section 3 — The JVM Analogy
% ══════════════════════════════════════════════════════════════════

\section{The JVM Analogy}\label{sec:analogy}

In 1995, Java promised \emph{``Write once, run anywhere.''}  The Java
Virtual Machine decoupled application code from the underlying hardware:
a developer could compile once and deploy on any platform that hosted
a conforming JVM.  The insight was that \textbf{execution} need not be
bound to a specific processor architecture.

NEOS makes an analogous promise one abstraction level higher:
\emph{``Think once, reason anywhere.''}  NEOS decouples
\textbf{reasoning} from any specific large language model.  Whether the
substrate is GPT-4, Claude, Gemini, Llama, or a model that does not yet
exist, the field dynamics, attractor basins, and collapse operators
behave identically---just as a Java program behaves identically on x86
and ARM.

The parallel is not merely rhetorical.  Table~\ref{tab:jvm-neos}
enumerates six structural correspondences between the JVM runtime and
the NEOS runtime, each grounded in a shared engineering insight.

% ── JVM vs NEOS table ────────────────────────────────────────────
\begin{table}[htbp]
\centering
\caption{Structural correspondences between the JVM and NEOS runtimes.}
\label{tab:jvm-neos}
\begin{tabularx}{\textwidth}{@{}llX@{}}
\toprule
\thead{JVM World} & \thead{NEOS World} & \thead{The Insight} \\
\midrule
Garbage Collection  & Decay ($\lambda$)
  & Automatic cleanup of unreferenced objects\,/\,unreinforced patterns \\
Heap Memory         & Field State
  & Where objects\,/\,patterns live and interact \\
Threads             & Multi-Field
  & Concurrent execution contexts with shared state \\
Stack Trace         & Cycle Trace
  & Debugging by tracing execution\,/\,dynamics history \\
ClassLoader         & Pattern Injection
  & Loading new code\,/\,meaning into the runtime \\
JIT Compilation     & Adaptive Resonance
  & Runtime optimisation based on actual usage patterns \\
\bottomrule
\end{tabularx}
\end{table}

\paragraph{The key argument.}
The JVM abstracted \emph{hardware}; NEOS abstracts \emph{cognition
itself}.  The JVM executes \textbf{code}---deterministic instruction
sequences that transform data.  NEOS executes \textbf{meaning}---patterns
that interact through resonance, compete for activation energy, and
self-organise into attractor basins.  Code is brittle: change one
instruction and the program crashes.  Meaning is robust: perturb a
pattern and the field relaxes back toward its attractors, provided the
perturbation falls within the basin of attraction.

\paragraph{The garbage-collection parallel.}
The correspondence between garbage collection and decay~($\lambda$)
deserves special attention.  In the JVM, the garbage collector reclaims
heap memory occupied by objects that are no longer reachable from any
live reference.  In NEOS, the decay term
$-\lambda\,A(x)$ continuously diminishes the activation of every pattern
that is not sustained by resonance with its neighbours.  Both mechanisms
solve the same fundamental problem---\emph{resource accumulation}---at
different abstraction levels: the JVM prevents memory leaks, NEOS
prevents \emph{semantic leaks}, the unbounded growth of irrelevant ideas
that would otherwise drown out meaningful signal.

% ── TikZ: JVM vs NEOS parallel stacks ───────────────────────────
\begin{figure}[htbp]
\centering
\begin{tikzpicture}[
  box/.style={
    draw, rounded corners=6pt, minimum width=3.8cm, minimum height=0.9cm,
    font=\small\sffamily, align=center, thick,
  },
  jvmbox/.style={box, fill=neosblue!8, draw=neosblue!40},
  neosbox/.style={box, fill=neosorange!8, draw=neosorange!50},
  conn/.style={dashed, gray!60, thick},
  lbl/.style={font=\footnotesize\sffamily, gray!70},
]

% Column titles
\node[font=\bfseries\sffamily, neosblue!60!black] at (-3.2, 5.8) {JVM Runtime};
\node[font=\bfseries\sffamily, neosorange!60!black] at ( 3.2, 5.8) {NEOS Runtime};

% Left column — JVM (top to bottom)
\node[jvmbox] (jit)  at (-3.2, 4.8) {JIT Compilation};
\node[jvmbox] (thr)  at (-3.2, 3.6) {Threads};
\node[jvmbox] (heap) at (-3.2, 2.4) {Heap Memory};
\node[jvmbox] (gc)   at (-3.2, 1.2) {Garbage Collection};
\node[jvmbox] (cl)   at (-3.2, 0.0) {ClassLoader};

% Right column — NEOS (top to bottom)
\node[neosbox] (ar)   at (3.2, 4.8) {Adaptive Resonance};
\node[neosbox] (mf)   at (3.2, 3.6) {Multi-Field};
\node[neosbox] (fs)   at (3.2, 2.4) {Field State};
\node[neosbox] (dec)  at (3.2, 1.2) {Decay ($\lambda$)};
\node[neosbox] (pi)   at (3.2, 0.0) {Pattern Injection};

% Dashed connections with centered arrows
\foreach \l/\r in {jit/ar, thr/mf, heap/fs, gc/dec, cl/pi} {
  \draw[conn] (\l.east) -- (\r.west)
    node[midway, font=\small] {$\leftrightarrow$};
}

\end{tikzpicture}
\caption{Parallel runtime stacks.  Each layer of the JVM has a
  structural correspondent in NEOS, operating on meaning rather than
  data.}
\label{fig:jvm-neos-stacks}
\end{figure}

Just as the JVM freed developers from thinking about registers and
calling conventions, NEOS frees analysts from thinking about prompt
engineering and model-specific idiosyncrasies.  The abstraction
boundary is the field equation: everything above it speaks the
language of patterns and resonance; everything below it is an
implementation detail of the substrate.
